\section{Tradução Automática Neural - NMT}
Como apresentado na Seção \ref{chap:NMT}, a tradução automática neural é uma
importante tarefa que tem como objetivo realizar a tradução de uma língua de
origem para uma língua de destino por meio de sistemas computacionais, mais
especificamente utilizando redes neurais artificiais. Essas redes buscam
simular, de forma simplificada, o funcionamento do cérebro humano, permitindo a
geração de traduções que se aproximem, o máximo possível, de traduções
realizadas por humanos \cite{tan2020neural}.

As seções seguintes têm como objetivo apresentar os conceitos teóricos que
fundamentam o estudo da arquitetura Transformer, proposta com o intuito de
aprimorar a tarefa de tradução automática. Essa arquitetura introduz novos
componentes e combina técnicas já existentes de forma mais eficiente. Como
resultado, a arquitetura Transformer estabeleceu um novo estado da arte na
tradução automática baseada em redes neurais artificiais
\cite{DBLP:journals/corr/VaswaniSPUJGKP17}.

\section{Tokenização}

Tokenização, ou \textit{tokenization} em inglês, é uma etapa de
pré-processamento fundamental no processamento de linguagem natural. Esse
processo consiste em converter textos em unidades menores, como bytes,
caracteres, subpalavras, palavras ou expressões de múltiplas palavras,
denominadas \textit{tokens} \cite{asgari2025morphbpemorphoawaretokenizerbridging}.

A Figura \ref{fig:tokenizer-rt} ilustra esse processo, no qual um texto é segmentado em tokens
e, posteriormente, cada token é representado por um numero, comumente chamado
de identificador ou id, cujo objetivo é mapear um token de volta à uma palavra
ao final de uma tradução entre línguas
\cite{asgari2025morphbpemorphoawaretokenizerbridging}.

\begin{figure}[H]
    \centering
    \includegraphics[width=.70\textwidth]{Imagens/tokenizer-referncial-teorico_compressed.pdf}

    \caption{\label{fig:tokenizer-rt}Tokenização de uma frase em português}

    \legend{Fonte: Criada pelo autor baseado na ferramenta \href{https://platform.openai.com/tokenizer}{openai tokenizer}}
\end{figure}

\section{Embedding}

O processamento de linguagem natural, em conjunto com tarefas de tradução entre
línguas, exige que palavras ou tokens sejam representados de forma numérica.
Nesse contexto, \textit{embeddings} correspondem ao processo de representação
de palavras ou tokens em um espaço vetorial de múltiplas dimensões, no qual
relações semânticas podem ser modeladas por meio de relações matemáticas \cite{llm_embeddings_explained}.

A representação de tokens em vetores de múltiplas dimensões é fundamental, pois
esses vetores são capazes de capturar informações semânticas relevantes. Dessa
forma, palavras com significados semelhantes tendem a possuir representações
vetoriais próximas nesse espaço. Por exemplo, os vetores associados às palavras
“gato” e “cachorro” tendem a ser mais próximos entre si do que em relação à
palavra “morango” \cite{llm_embeddings_explained}.

Essa propriedade é essencial para a geração de traduções mais naturais e
coerentes, uma vez que permite aos modelos capturar relações semânticas entre
palavras e contextos linguísticos \cite{llm_embeddings_explained}.

A Figura \ref{fig:embedding-rt} ilustra o processo de conversão de uma sequência de tokens em
representações vetoriais em um espaço de multiplas dimensões.

\begin{figure}[H]
    \centering
    \includegraphics[width=.90\textwidth]{Imagens/embedding-rt.pdf}

    \caption{\label{fig:embedding-rt}Conversão de tokens em embeddings}

    \legend{Fonte: \cite{llm_embeddings_explained} adaptada pelo autor}
\end{figure}

\section{Função Softmax}

\textit{Softmax} é uma função matemática frequentemente utilizada como função de ligação probabilística para dados categóricos não ordenados. Dado um vetor de entradas, a aplicação da função softmax resulta em uma saída composta por uma distribuição de probabilidades, na qual cada valor está associado a um elemento da entrada \cite{franke2023softmax}.

A função softmax é amplamente utilizada em modelos de aprendizado de máquina,
especialmente em tarefas que envolvem classificação e tradução automática,
sendo responsável por converter pontuações geradas pelo modelo em distribuições
de probabilidade sobre os possíveis resultados de saída \cite{franke2023softmax}.

A função softmax é definida na equação \ref{eq:softmax-definition}:
\begin{equation}
    p_i = \frac{\exp(\alpha s_i)}{\sum_j \exp(\alpha s_j)}
    \label{eq:softmax-definition}
\end{equation}

onde:

\begin{itemize}
    \item \(\alpha\): parâmetro de escala da função softmax, frequentemente omitido e considerado igual a 1;
    \item \(s_i\): valor da i-ésima posição do vetor de pontuações.
\end{itemize}

A Figura \ref{fig:softmax1-rt} ilustra o funcionamento da função softmax
aplicada a um vetor de pontuações, resultando na geração de uma distribuição de
probabilidades associada a cada elemento de entrada.

\begin{figure}[H]
    \centering
    \includegraphics[width=.90\textwidth]{Imagens/softmax-1-rt_compressed.pdf}
    \caption{\label{fig:softmax1-rt}Conversão de pontuações em probabilidades}
    \legend{Fonte: \cite{franke2023softmax} adaptada pelo autor}
\end{figure}

Uma das propriedades mais importantes da função softmax em modelos neurais é a
sua capacidade de normalização, garantindo que a soma das probabilidades
geradas seja igual a 1. Essa característica permite interpretar a saída do
modelo como uma distribuição de probabilidade válida sobre os possíveis
resultados. A Figura \ref{fig:softmax2-rt} ilustra essa propriedade
\cite{franke2023softmax}.

\begin{figure}[H]
    \centering
    \includegraphics[width=.70\textwidth]{Imagens/softmax2-rt_compressed.pdf}
    \caption{\label{fig:softmax2-rt}Propriedade de normalização da função softmax}
    \legend{Fonte: \cite{franke2023softmax}}
\end{figure}

\section{Mecanismo De Atenção}

O mecanismo de atenção é uma técnica cujo objetivo é identificar e relacionar
as partes mais relevantes dos dados de entrada durante o processamento de
sequências \cite{hays2026attentionmechanismsneuralnetworks}.

Grande parte das arquiteturas de tradução automática baseadas em redes neurais
utiliza dois componentes principais: o codificador (encoder) e o decodificador
(decoder), conforme apresentado na Seção \ref{chap:NMT}. No entanto, modelos
tradicionais baseados apenas nessa estrutura enfrentam limitações ao
representar sequências longas, uma vez que toda a informação precisa ser
comprimida em um vetor de dimensão fixa \cite{hays2026attentionmechanismsneuralnetworks}.

Nesse contexto, os mecanismos de atenção modificam a forma como redes neurais
processam dados sequenciais, permitindo acesso direto a todas as posições da
sequência de entrada. Isso elimina a necessidade de compressão em uma única
representação vetorial e possibilita uma melhor captura de relações
contextuais \cite{hays2026attentionmechanismsneuralnetworks}.

A operação de atenção consiste na computação de uma soma ponderada sobre os
dados de entrada, na qual cada peso representa a importância de uma determinada
posição da sequência para o elemento que está sendo processado. Na prática,
isso permite que diferentes partes da entrada se relacionem entre si,
possibilitando a captura de dependências contextuais entre palavras em tarefas
de tradução \cite{hays2026attentionmechanismsneuralnetworks}.

As equações \ref{eq:attention-score}, \ref{eq:attention-softmax} e \ref{eq:attention-calculation} definem a formulação matemática de um mecanismo de atenção
\cite{hays2026attentionmechanismsneuralnetworks}:

\begin{equation}
    e_{ij} = score(\mathbf{q}_i, \mathbf{k}_j)
    \label{eq:attention-score}
\end{equation}

\begin{equation}
    \alpha_{ij} = softmax(e_{ij})
    \label{eq:attention-softmax}
\end{equation}

\begin{equation}
    \mathbf{c}_i = \sum_{j=1}^{n} \alpha_{ij} \mathbf{v}_j
    \label{eq:attention-calculation}
\end{equation}

Onde:

\begin{itemize}
    \item \(\mathbf{q}_i\): vetor de consulta (\textit{query}), representando a informação buscada na posição \(i\);
    \item \(\mathbf{k}_j\): vetor de chave (\textit{key}), representando a informação disponível na posição \(j\);
    \item \(\mathbf{v}_j\): vetor de valor (\textit{value}), contendo a informação a ser agregada;
    \item \(\alpha_{ij}\): peso de atenção, formando uma distribuição de probabilidade sobre as posições de entrada, indicando o grau de contribuição de cada posição para a saída na posição \(i\).
\end{itemize}

A Figura \ref{fig:attention-1-rt} ilustra um exemplo de tradução automática sem
o uso de mecanismos de atenção, no qual é possível observar limitações na
captura de relações entre palavras da língua de origem e da língua de destino.

\begin{figure}[H]
    \centering
    \includegraphics[width=.70\textwidth]{Imagens/attention-rt-1_compressed.pdf}
    \caption{\label{fig:attention-1-rt}Tradução sem a utilização de um mecanismo de atenção}
    \legend{Fonte: \cite{raschka2023selfattention}}
\end{figure}

Por outro lado, a Figura \ref{fig:attention-2-rt} apresenta o uso do mecanismo
de atenção, evidenciando a capacidade do modelo de capturar relações entre
palavras, resultando em traduções mais coerentes e contextualizadas.

\begin{figure}[H]
    \centering
    \includegraphics[width=.70\textwidth]{Imagens/attention-rt-2_compressed.pdf}
    \caption{\label{fig:attention-2-rt}Tradução com a utilização de um mecanismo de atenção}
    \legend{Fonte: \cite{raschka2023selfattention}}
\end{figure}

\section{Arquitetura Transformer}
\subsection{Visão Geral}

A arquitetura \textit{Transformer} foi proposta em 2017 pela empresa Google no artigo
intitulado \textit{“Attention Is All You Need”}, representando uma mudança de
paradigma em relação aos modelos até então predominantes, baseados em redes
neurais recorrentes (RNN).

Essa arquitetura introduziu avanços importantes, como a capacidade de capturar
dependências de longo alcance entre elementos de uma sequência e o
processamento paralelo dos dados, resultando em um novo estado da arte em
diversas tarefas de processamento de linguagem natural.

A partir da arquitetura Transformer, tornou-se possível o desenvolvimento de
modelos avançados de geração de texto, como o GPT, da OpenAI, o Gemini, da
Google, e o LLaMA, da Meta, além de impulsionar a evolução de modelos de
tradução automática, como o MarianMT, de código aberto, e os modelos M2M e
NLLB, também desenvolvidos pela Meta. Esses modelos representam avanços
significativos na tarefa de tradução entre línguas, todos baseados nos
princípios introduzidos pela arquitetura Transformer
\cite{build-llms-from-scratch-book}.

\begin{figure}[H]
    \centering
    \includegraphics[height=0.6\textwidth]{Imagens/transformer_compressed.pdf}
    \caption{\label{fig:fig-architecture}Arquitetura Transformer}
    \legend{Fonte: \cite{DBLP:journals/corr/VaswaniSPUJGKP17}}
\end{figure}
\subsection{Positional Encoding}
O \textit{positional encoding}, ou codificação posicional, é uma técnica
utilizada para incorporar informações sobre a posição dos \textit{tokens} em uma
sequência na arquitetura Transformer.

Em tarefas de tradução automática, a ordem das palavras desempenha um papel
fundamental na construção do significado semântico de uma sentença. A posição
dos termos permite identificar relações de dependência entre palavras, além de
influenciar diretamente a estrutura da frase na língua de destino, uma vez que
diferentes idiomas apresentam diferentes organizações sintáticas.

Como a arquitetura Transformer não possui mecanismos recorrentes ou
convolucionais, ela não é capaz de capturar informações de ordem de forma
implícita. Para contornar essa limitação, são adicionadas informações
posicionais aos embeddings de entrada por meio de funções determinísticas
baseadas em seno e cosseno \cite{ibm_positional_encoding}.

A codificação posicional é definida pelas equações \ref{eq:positional-encoding-sin} e \ref{eq:positional-encoding-cos}
\cite{DBLP:journals/corr/VaswaniSPUJGKP17}:

\begin{equation}
    PE_{(pos, 2i)} = \sin\left(\frac{pos}{10000^{2i/d_{model}}}\right)
    \label{eq:positional-encoding-sin}
\end{equation}

\begin{equation}
    PE_{(pos, 2i+1)} = \cos\left(\frac{pos}{10000^{2i/d_{model}}}\right)
    \label{eq:positional-encoding-cos}
\end{equation}

Onde:

\begin{itemize}
    \item \( pos \): posição do token na sequência;
    \item \( i \): índice da dimensão do vetor de embedding;
    \item \( d_{model} \): dimensão do vetor de embedding;
    \item \( PE \): função de codificação posicional que associa uma posição da sequência a uma representação vetorial.
\end{itemize}
Essas funções permitem que cada posição na sequência seja representada por um vetor único, preservando relações relativas entre posições, o que auxilia o modelo na identificação da ordem dos \textit{tokens}.

A Figura \ref{fig:fig-positiona-encoding} ilustra o processo de aplicação do
\textit{positional encoding} em uma sequência de entrada.

\begin{figure}[H]
    \centering
    \includegraphics[height=0.6\textwidth]{Imagens/positional-encoding_compressed.pdf}
    \caption{\label{fig:fig-positiona-encoding} Exemplo de codificação posicional aplicada a uma sequência}
    \legend{Fonte: \cite{brownlee_positional_encoding} adaptado}
\end{figure}

\subsection{Feed-Forward Neural Network}
Redes neurais \textit{feed-forward} são modelos computacionais nos quais a
informação flui em uma única direção, da entrada para a saída, sem a presença
de ciclos ou realimentação. Essas redes são compostas por camadas de neurônios
interconectados, nas quais cada conexão possui um peso associado, responsável
por ajustar a influência dos dados ao longo do processamento \cite{sazli2006brief}.

De forma geral, uma rede \textit{feed-forward} é estruturada em três tipos de
camadas: camada de entrada, responsável por receber os dados, camadas ocultas,
responsáveis pelo processamento das informações por meio de transformações não
lineares, e camada de saída, que produz o resultado final do
modelo \cite{sazli2006brief}.

No contexto da arquitetura Transformer, as redes \textit{feed-forward} são
aplicadas de forma independente a cada posição da sequência, após o mecanismo
de atenção. Seu principal objetivo é transformar e refinar as representações
vetoriais geradas pelas camadas de atenção, contribuindo para a modelagem de
relações mais complexas entre os elementos da
sequência.\cite{gerber2025attentionneedimportancefeedforward}

A Figura \ref{fig:fig-feed-forward} ilustra um exemplo de rede neural
\textit{feed-forward} com uma camada oculta.

\begin{figure}[H]
    \centering
    \includegraphics[width=0.6\textwidth]{Imagens/feed-forward_compressed.pdf}
    \caption{\label{fig:fig-feed-forward} Rede neural feed-forward}
    \legend{Fonte: \cite{geeksforgeeks_fnn} adaptado}
\end{figure}

\subsection{Scaled Dot-Product Attention}
No contexto da arquitetura Transformer, o mecanismo de atenção é implementado
por meio da chamada \textit{Scaled Dot-Product Attention}. Esse mecanismo
calcula a relevância entre diferentes elementos da sequência de entrada
utilizando três componentes principais: consultas (\textit{queries}), chaves
(\textit{keys}) e valores (\textit{values}).

A atenção é calculada por meio do produto escalar entre consultas e chaves,
seguido de uma normalização utilizando a função softmax, conforme a equação \ref{eq:attention-mecanism}
\cite{DBLP:journals/corr/VaswaniSPUJGKP17}:

\begin{equation}
    {Attention}(Q, K, V) = {softmax}\left( \frac{QK^T}{\sqrt{d_k}} \right)V
    \label{eq:attention-mecanism}
\end{equation}

onde:

\begin{itemize}
    \item \( Q \): Representa a perspectiva da palavra atual. Para uma palavra específica \( i \), seu vetor que consulta(\( Q \)) é utilizado para comparar com todas as palavras disponíveis;
    \item \( K \): Representa o rótulo ou identificador da informação que uma palavra carrega. Cada palavra \( j \) possui um vetor de chaves(\( K \)), onde cada consulta \( Q_i \) é comparada com todas as chaves, \( K_i \), com o objetivo do calcular a atenção entre as palavras, indicando o quão relevante a palavra \( j \) é para a palavra \( i \);
    \item \( V \): Representa o conteúdo atual ou significado de uma palavra atual. Onde cada palavra \( j \) possui também um vetor de valores \( v_j \). Uma vez que que a atenção entre \( Q_i \) e \( K_j \) é computado, este valor é utilizado para ajustar a relevância do vetor de valores correspondentes \( V_j \) \cite{ibm_attention_mechanism}.
\end{itemize}

O fator de escala \(\sqrt{d_k}\) é utilizado para evitar valores muito elevados
no produto escalar, o que poderia prejudicar a estabilidade da função softmax.

\subsection{Multi-Head Attention}
O mecanismo de \textit{multi-head attention} é uma extensão do mecanismo de
atenção utilizado na arquitetura Transformer, cujo objetivo é permitir que o
modelo capture diferentes tipos de relações entre os elementos de uma sequência
de forma paralela.

Diferentemente da atenção simples, que realiza uma única projeção sobre as
matrizes de consultas, chaves e valores, o \textit{multi-head attention}
projeta essas matrizes múltiplas vezes em diferentes subespaços vetoriais. Cada
uma dessas projeções é processada por um mecanismo de atenção independente,
denominado \textit{head}.

Essa abordagem permite que o modelo aprenda diferentes representações e
relações entre os \textit{tokens}, como dependências sintáticas e semânticas,
contribuindo para uma compreensão mais rica do contexto da sequência
\cite{liu2021multi}.

A formulação matemática do mecanismo de \textit{multi-head attention} é dada
pelas equações \ref{eq:multi-head-attention} e  \ref{eq:head-mecanism-attention} \cite{DBLP:journals/corr/VaswaniSPUJGKP17}:

\begin{equation}
    {MultiHead}(Q, K, V) = {Concat}({head}_1, \dots, {head}_h)W^O
    \label{eq:multi-head-attention}
\end{equation}

onde cada \textit{head} é definido como:

\begin{equation}
    {head}_i = {Attention}(QW_i^Q, KW_i^K, VW_i^V)
    \label{eq:head-mecanism-attention}
\end{equation}

Nesse contexto, \(W_i^Q\), \(W_i^K\) e \(W_i^V\) são matrizes de projeção
aprendidas para cada \textit{head} de atenção, enquanto \(W^O\) é uma matriz
responsável por combinar as saídas concatenadas das diferentes \textit{heads}
\cite{DBLP:journals/corr/abs-2006-16362}.

\subsection{Encoder}
Encoder, ou codificador, é um dos principais componentes da arquitetura
Transformer, responsável por receber a sequência de entrada já convertida em
embeddings e enriquecida com informações de posição por meio do
\textit{positional encoding}. Seu objetivo é transformar essa sequência em um
conjunto de representações vetoriais contextuais, capazes de capturar relações
semânticas relevantes entre os \textit{tokens}.

Cada bloco do encoder é composto por duas subcamadas principais. A primeira é o
mecanismo de atenção (\textit{self-attention}), responsável por relacionar os
\textit{tokens} da sequência entre si, permitindo a captura de dependências de curto e
longo alcance. A segunda subcamada é a rede \textit{feed-forward}, aplicada de
forma independente a cada posição da sequência, com o objetivo de transformar e
enriquecer as representações vetoriais geradas pela etapa anterior.

Após cada subcamada, são aplicadas operações de normalização e conexões
residuais, com a finalidade de estabilizar o treinamento e preservar
informações importantes ao longo das camadas do modelo.

Ao final do processamento pelo encoder, é gerado um conjunto de vetores de
representações contextuais, no qual cada vetor está associado a um token da
sequência de entrada. Esses vetores incorporam informações sobre o contexto
global da sentença, incluindo relações entre palavras, e são posteriormente
utilizados pelo decoder para a geração da tradução na língua de destino
\cite{build-llms-from-scratch-book}.

\subsection{Decoder}
Decoder, ou decodificador, é o componente da arquitetura Transformer
responsável por gerar a sequência de saída, ou seja, a tradução na língua de
destino, com base nas representações contextuais produzidas pelo encoder.

Assim como o encoder, o decoder é composto por múltiplos blocos empilhados,
sendo cada bloco formado por três subcamadas principais: \textit{Masked
    Self-Attention}, \textit{Cross-Attention} (atenção sobre o encoder) e uma rede
\textit{feed-forward}.

A subcamada de \textit{Masked Self-Attention} aplica o mecanismo de atenção
sobre a própria sequência de saída, porém com uma máscara que impede o modelo
de acessar \textit{tokens} futuros. Dessa forma, a geração ocorre de maneira
autoregressiva, considerando apenas os \textit{tokens} já produzidos até o momento, o
que garante coerência na construção da tradução.

A subcamada de \textit{Cross-Attention} permite que o decoder utilize as
informações produzidas pelo encoder, estabelecendo relações entre os \textit{tokens} da
sequência de saída e os \textit{tokens} da sequência de entrada. Esse mecanismo é
essencial para alinhar corretamente a tradução com o contexto da frase
original.

Além disso, o mecanismo de atenção utilizado é do tipo \textit{multi-head
    attention}, que realiza múltiplas operações de atenção em paralelo, permitindo
ao modelo capturar diferentes tipos de relações entre os elementos da
sequência.

Por fim, a camada \textit{feed-forward} é aplicada com o objetivo de
transformar e refinar as representações geradas, contribuindo para a produção
de uma saída mais consistente e contextualizada.

Durante o processo de geração, o decoder produz a tradução de forma iterativa,
gerando um token por vez. A cada etapa, o token gerado é realimentado como
entrada para o próprio decoder, juntamente com as representações do encoder,
até que um token especial de término seja produzido, indicando o fim da
tradução \cite{build-llms-from-scratch-book}.

\section{Modelos de Tradução Open Source}
\subsection{MarianMT}
MarianMT é um \textit{framework} de tradução automática neural desenvolvido por
pesquisadores da Adam Mickiewicz University, em Poznań, e da University of
Edinburgh. Implementado majoritariamente na linguagem C++, o MarianMT foi
projetado com foco em eficiência computacional e baixo número de dependências
externas.

Diferentemente de modelos prontos para uso, o MarianMT consiste em uma
plataforma para desenvolvimento, treinamento e implementação de modelos de
tradução automática neural. Ele oferece suporte à construção de modelos
baseados em diferentes arquiteturas, incluindo redes neurais recorrentes e,
principalmente, modelos baseados na arquitetura Transformer.

Uma das principais características do MarianMT é sua otimização para
desempenho, permitindo o treinamento eficiente em múltiplas unidades de
processamento, como CPUs e GPUs. Além disso, o framework possibilita a criação
de modelos especializados para diferentes pares de línguas, a partir de bases
de dados específicas \cite{mariannmt}.

Entre os principais aspectos do MarianMT, destacam-se \cite{mariannmt}:

\begin{itemize}
    \item Implementação eficiente em C++, com foco em desempenho;
    \item Suporte a treinamento e inferência em CPU e GPU;
    \item Compatibilidade com arquiteturas modernas de tradução neural, como Transformer;
    \item Código aberto sob licença permissiva (MIT).
\end{itemize}

\subsection{M2M-100}
\label{section-M2M}
M2M-100 (\textit{Many-to-Many}) é um modelo de tradução automática neural baseado na arquitetura transfomer e desenvolvido pela Meta, cujo objetivo é realizar traduções diretas entre múltiplos pares de línguas, mais precisamente entre 100 línguas, sem a necessidade de utilizar o inglês como língua intermediária.

Em abordagens tradicionais, quando a língua de origem ou destino não envolve o
inglês, é comum a utilização de uma estratégia intermediária, na qual a
tradução é realizada inicialmente para o inglês e, posteriormente, convertida
para a língua de destino. Essa abordagem é frequentemente adotada devido à
maior disponibilidade de dados de treinamento em inglês.

O modelo M2M-100 foi proposto para superar essa limitação, permitindo traduções
diretas entre diferentes línguas. Para isso, o modelo foi treinado em uma
grande base de dados multilíngue, construída com o objetivo de abranger uma
ampla variedade de pares linguísticos.

Além disso, foram utilizadas estratégias para melhorar a qualidade dos dados,
incluindo o agrupamento de línguas com base em similaridades linguísticas,
geográficas e culturais. Essas técnicas contribuíram para a geração de
traduções mais consistentes e contextualizadas, especialmente para pares de
línguas com menor disponibilidade de dados.

O desenvolvimento do M2M-100 atende à necessidade de plataformas globais, como
as aplicações da Meta, de oferecer suporte à comunicação entre usuários de
diferentes idiomas com traduções de alta qualidade \cite{fan2020m2m100}.

\subsection{NLLB}
NLLB (\textit{No Language Left Behind}) é um modelo de tradução automática
neural baseado na arquitetura Transformer, desenvolvido pela Meta, com o
objetivo de realizar traduções de alta qualidade entre um grande número de
línguas, incluindo línguas com poucos recursos disponíveis.

Diferentemente do modelo M2M-100, apresentado na Seção \ref{section-M2M}, o
NLLB representa uma evolução significativa, sendo projetado para suportar mais
de 200 línguas, com foco especial na melhoria da qualidade de tradução para
línguas de baixo recurso (\textit{low-resource languages}).

Um dos principais desafios enfrentados no desenvolvimento do NLLB está
relacionado à escassez de dados paralelos para determinadas línguas. Para
contornar esse problema, foram utilizadas estratégias de construção e curadoria
de bases de dados multilíngues em larga escala, incluindo a coleta de textos de
diversas fontes, como entrevistas com nativos da língua e o uso de técnicas de
filtragem para garantir maior qualidade dos dados.

Além disso, o treinamento do modelo envolveu o uso de técnicas avançadas de
otimização e regularização, com o objetivo de melhorar a generalização e
reduzir problemas como o \textit{overfitting}, especialmente em cenários com
dados limitados \cite{nllbteam2022languageleftbehindscaling}

\section{Métricas de Avaliação}
\subsection{BLEU}


A métrica BLEU (\textit{Bilingual Evaluation Understudy}) é um método automático utilizado para avaliação de traduções automáticas. Trata-se de uma métrica independente de linguagem, cujo objetivo é medir o grau de similaridade entre uma tradução gerada automaticamente e uma ou mais traduções de referência produzidas por humanos.

A ideia central da métrica BLEU é que traduções automáticas de maior qualidade tendem a apresentar maior sobreposição com traduções humanas de referência. Dessa forma, quanto maior a similaridade entre a tradução gerada e a tradução de referência, maior será o valor da métrica BLEU.

O valor da métrica BLEU varia no intervalo de 0 a 1, onde valores mais próximos de 1 indicam maior similaridade com a tradução de referência, enquanto valores próximos de 0 indicam baixa similaridade \cite{papineni2002bleu}.

A formulação matemática da métrica BLEU é definida como \ref{eq:bleu}

\begin{equation} \label{eq:bleu}
\begin{aligned}
    \text{BP}   & =
    \begin{cases}
        1           & \text{if } c > r    \\
        e^{(1-r/c)} & \text{if } c \leq r
    \end{cases},  \\[6pt]
    \text{BLEU} & = \text{BP} \cdot \exp \left( \sum_{n=1}^{N} w_n \log p_n \right) .
\end{aligned}
\end{equation}

Onde:

\begin{itemize}
    \item \(c\): comprimento da tradução gerada;
    \item \(r\): comprimento da tradução de referência;
    \item \(w_n\): pesos associados a cada ordem de n-grama, tais que \( \sum w_n = 1 \);
    \item \(p_n\): precisão modificada de n-gramas.
\end{itemize}

\subsection{BERTScore}
    BERTScore é uma métrica automática utilizada para avaliação de textos gerados por modelos de processamento de linguagem natural, incluindo tarefas de tradução automática.

    Diferentemente de métricas baseadas em correspondência exata de n-gramas, como o BLEU, o BERTScore avalia a similaridade entre uma sentença gerada e uma sentença de referência por meio de representações vetoriais contextuais (\textit{embeddings}) obtidas a partir de modelos pré-treinados, como o BERT.

    O cálculo do BERTScore é realizado por meio da comparação entre os embeddings de cada token da sentença gerada e os embeddings dos \textit{tokens} da sentença de referência. Para cada token, é calculada uma medida de similaridade (geralmente baseada em cosseno), permitindo capturar relações semânticas mesmo quando não há correspondência exata entre as palavras.

    A avaliação por meio do BERTScore é composta pelas métricas de precisão, revocação e F1-score.

    Onde:

\begin{itemize}
    \item Precisão: considera a sentença gerada pelo modelo e mede o quanto os \textit{tokens} gerados estão semanticamente alinhados com a sentença de referência;
    
    \item Revocação: considera a sentença de referência e avalia o quanto seu conteúdo está presente na sentença gerada pelo modelo;
    
    \item F1-score: combina precisão e revocação em um único valor, representando um equilíbrio entre a cobertura do conteúdo da referência e a qualidade da correspondência semântica.
\end{itemize}

    O valor do F1-score é frequentemente utilizado como resultado principal da métrica BERTScore, sendo que valores mais elevados indicam maior similaridade semântica entre a sentença gerada e a sentença de referência \cite{zhang2019bertscore}.


\subsection{chrF}
    ChrF (\textit{CHaracter n-gram F-score}) é uma técnica utilizada para avaliação automática de modelos de tradução, baseada na comparação de n-gramas de caracteres entre a sentença gerada e a sentença de referência \cite{popovic2015chrf}.

    Diferentemente de métricas como o BLEU, que operam no nível de palavras, o chrF realiza a comparação no nível de caracteres, permitindo maior sensibilidade a variações morfológicas e flexionais presentes nas línguas naturais.

    O cálculo da métrica chrF baseia-se na combinação de precisão e revocação de n-gramas de caracteres, resultando em uma medida do tipo F-score. Essa abordagem permite capturar similaridades entre palavras mesmo quando há pequenas variações, como flexões ou erros parciais na tradução 
    \cite{emergentmind2026chrf}.


    A formulação matemática da métrica chrF é definida como \ref{eq:chrf-formula} \cite{popovic2015chrf}:
    
    \begin{equation}
    \text{CHRF}\beta = (1 + \beta^2) \frac{\text{CHRP} \cdot \text{CHRR}}{\beta^2 \cdot \text{CHRP} + \text{CHRR}}
    \label{eq:chrf-formula}
    \end{equation}

onde:

\begin{itemize}
    \item \(\text{chrP}\): precisão de n-gramas de caracteres, representando a proporção de n-gramas presentes na sentença gerada que também aparecem na referência;
    
    \item \(\text{chrR}\): revocação de n-gramas de caracteres, representando a proporção de n-gramas da referência que são recuperados pela sentença gerada;
    
    \item \(\beta\): parâmetro que controla a importância relativa entre precisão e revocação. Quando \(\beta = 1\), ambas possuem o mesmo peso.
\end{itemize}




